% ============================================================
% PAPER 4 of 11 -- minimax_dual_certificates
% Assembled 2026-09-29 by content-and-merits partition.
% Title: Minimax Dual Certificates and the Obstruction Taxonomy
% Source: arena agent 1/agent workspace/fam/minimax_v11.tex
% Partition note: Single-source. Duality for the obstruction programme, with its own two-stock benchmark.
%% STATUS: structurally assembled as a NEW version. Editorial integration
%%         (de-duplication of shared framework, sibling cross-citations)
%%         was reviewed for this editorial version; see the editorial change log.
%% ============================================================

%% v14 (2026-09-29): added section \ref{priorart} and seven bibliography entries, all
%%   VERIFIED against publisher records before use.
%%   HONEST BASELINE: this paper's obstruction IS THE EMPTINESS OF THE DISCRIMINATING
%%   KERNEL. That collision with Aubin / Aubin-Catte / Cardaliaguet-Quincampoix-Saint-
%%   Pierre is conceded in the text, not argued around. What survives is the CERTIFICATE
%%   FORM (a finitely supported MEASURE with a tight k+1 bound, rather than a subset of
%%   state space characterised by a fixed point) and the NEGATIVE BOUNDARY RESULT
%%   (prop:gap), which is the complement of CQS 1994 Prop 2.3: convexity is exactly what
%%   makes kernels convex AND what makes measure duality sound.
%%   CITATION NOTES: Aubin and CATTE (accented). The 1994 RAIRO/M2AN paper is the
%%   discriminating-kernel ALGORITHM paper. Sion 1958 supplies the minimax EQUALITY but
%%   not the sparse witness -- that distinction is drawn explicitly.
%%   DROPPED: the 'Aubin 2001, Poincare's shadow' attribution from the earlier roadmap
%%   could NOT be verified and is NOT asserted. Aubin 2001 is cited only for the
%%   verified kernel/capture-basin fixed-point characterisation.
%% PRIOR ART REQUIRED -- bar item, inserted 2026-09-29.
%% No related-work position. It must engage:
%%   - classical Farkas / linear-programming duality as the certificate mechanism,
%%     and what is genuinely new over the classical alternative.
%%   - Cardaliaguet, Quincampoix and Saint-Pierre (2007): discriminating kernels
%%     in the differential-game formulation.
%%   - Aubin and Catte (2002): bilateral fixed-point and discriminating-kernel
%%     calculus.
%%   - Aubin (2001): the complement of the kernel -- Poincare's "shadow".
%% Claim to establish: these dual certificates are the one-sided,
%% observation-constrained analogue of the discriminating-kernel programme, and
%% state which direction (necessity or sufficiency) each existing result serves.

% SIBLING PAPERS: Paper 1 (Obstruction Certificates under Incomplete Observation), Paper 2 (Exact Probabilistic Sufficiency), Paper 3 (Computational Certification of Obstruction), Paper 5 (Exact Belief Computation).
% Cross-cite, do not re-derive: the selector principle, the epistemic kernel
% and the observation structure are defined in Paper 1. Papers 2-5 cite it.

% SIBLING PAPERS: Paper 1 (Obstruction Certificates under Incomplete Observation), Paper 2 (Exact Probabilistic Sufficiency), Paper 3 (Computational Certification of Obstruction), Paper 5 (Exact Belief Computation).
% Cross-cite, do not re-derive: the selector principle, the epistemic kernel
% and the observation structure are defined in Paper 1. Papers 2-5 cite it.

\documentclass[10pt,twocolumn]{article}
\usepackage[a4paper,margin=15mm]{geometry}
\usepackage{amsmath,amssymb}
\usepackage{mathrsfs}
\usepackage{amsthm}
\theoremstyle{definition}
\newtheorem{theorem}{Theorem}
\newtheorem{proposition}{Proposition}
\newtheorem{corollary}{Corollary}
\newtheorem{remark}{Remark}
\newtheorem{definition}{Definition}
\newtheorem{conjecture}{Conjecture}
\newtheorem{lemma}{Lemma}
\usepackage{graphicx}
\usepackage{booktabs,array,calc}
\usepackage[font=small,labelfont=bf]{caption}
\usepackage[colorlinks=true]{hyperref}
\usepackage{etoolbox}
\providecommand{\tightlist}{\setlength{\itemsep}{0pt}\setlength{\parskip}{0pt}}
\setcounter{secnumdepth}{-1}
\emergencystretch=3em
\allowdisplaybreaks
\AtBeginDocument{%
  \BeforeBeginEnvironment{equation}{\small}%
  \BeforeBeginEnvironment{equation*}{\small}%
  \BeforeBeginEnvironment{align}{\small}%
  \BeforeBeginEnvironment{align*}{\small}%
  \BeforeBeginEnvironment{gather}{\small}%
  \BeforeBeginEnvironment{gather*}{\small}%
}

\begin{document}
\title{Minimax Dual Certificates for the\\ Common-Action Obstruction}
\author{Amin Abaee\\[0.35em]
{\small Independent Researcher, Tehran, Iran}\\[0.55em]
{\small\href{https://orcid.org/0000-0002-0019-1842}{ORCID: 0000-0002-0019-1842}}\\[0.3em]
{\small\href{mailto:amin\_abaee@ut.ac.ir}{amin\_abaee@ut.ac.ir}}}
\date{September 25, 2026}
\maketitle

\begin{abstract}
The common-action obstruction certifies that no single instrument serves
every branch of an information set; its polyhedral form carries a Farkas
multiplier (Abaee, 2026, An obstruction calculus, Theorem 2). This paper states and proves the
measure dual of the same obstruction: for control-affine dynamics with a
compact convex control set, the common safe-action set is empty precisely
when an adversarial probability measure on the active
boundary--disturbance bundle drives the expected inward drift strictly
negative against every control, with a certificate supported on at most
$k+1$ points ($k$ the control dimension; tightness is exhibited here in dimensions one and two). The
polyhedral case recovers the $l_1$-normalized Farkas pair exactly; the
singleton case recovers the Isaacs drift with the adversarial
\emph{distribution} over disturbances load-bearing. A two-action instance
exhibits a strict minimax gap: without convexity the equivalence fails and
no measure certifies --- the audited finite codex systems sit outside this
scope by design, certified combinatorially. An obstruction taxonomy
(physical, epistemic, institutional) is instantiated on the audited grid,
and a calibrated two-stock resource model --- one fixed control set
throughout --- yields a closed-form critical aggregate
$Y^{*} = \tfrac{27}{5}$: aggregate-yield fibres are viable exactly up to
$Y^{*}$, the obstruction at $Y = 6$ is certified by a two-point dual
measure with margin $\tfrac{3}{50}$, and the fibre-width formula
$\Delta x_{1}(Y) = Y - 4$ gives the aggregation rule its exact reading.
Dynamic-envelope, quantization, and stochastic-bridge extensions are
recorded as conjectures; the residual dynamic-observation completeness gap
of the source paper remains open. Eight exact-rational check families verify the declared finite
instances, not the general analytic converse or dynamic-envelope claims.
\end{abstract}

\vspace{0.4em}
\noindent\textbf{Keywords:} minimax duality; Farkas certificates; zero-sum
games; obstruction certificates; adversarial priors; Isaacs condition;
exact arithmetic

\section{Introduction}

The obstruction calculus's certificates are emptiness statements: a belief is
obstructed when a correspondence of jointly acceptable responses is empty
(Abaee, 2026, An obstruction calculus, Theorem 2; Abaee, 2026, Worked systems for the obstruction calculus). For polyhedral safe-action sets the
emptiness carries an explicit dual object --- the Farkas multiplier vector (Farkas, 1902)
--- so that the obstruction itself becomes a checkable artifact; the
sparse-witness bound (at most $m+1$ compatible states) and the singleton
Isaacs drift (Isaacs, 1965) complete the static picture. This paper unifies the static layer: the measure dual of the common-action obstruction, of which the normalized Farkas pair, the sparse witness, and the Isaacs reduction are special cases. The measure dual is verified exactly; the sparsity bound is proved by a finite-reduction route (finite reduction, linear-programming duality, basic-solution support); the singleton Isaacs reading is stated with its load-bearing distributional caveat; the equality's convexity boundary is exhibited by counterexample; the state-versus-epistemic separation is adopted as a working axiom; and a calibrated two-stock benchmark with a single fixed control set is solved in closed form. The dynamic layer --- an adversarial-prior envelope for the restricted-information game --- is then formulated and analyzed, and the residual dynamic-observation completeness gap (Abaee, 2026, Obstruction certificates under incomplete observation, Open Problem~1 in the revised manuscript) is \emph{not} closed here. The contributions are: the measure dual of the common-action obstruction, exactly verified, with the normalized Farkas pair, the sparse witness, and the Isaacs reduction as special cases, and the equality's convexity boundary exhibited by counterexample (Section 2); the calibrated two-stock benchmark solved in closed form (Section 4); and the dynamic layer --- the adversarial-prior envelope for the restricted-information game, its three obstructions, and the review-sequence envelope's exact couplings (Sections 5--6).

\begin{table}[h]
\centering\footnotesize
\caption{The theorem and its check families: what the verification record
certifies.}
\label{tab:checks}
\begin{tabular}{@{}l p{0.82\linewidth}@{}}
\toprule
Family & Check family / reported scope \\
\midrule
G1 & general tower on all trees ($K{=}3$, adaptive filtration); \\
   & martingale property at every cell, all 8192 policies \\
G2 & general DP equals path enumeration ($K{=}2$, 19683 policies); \\
   & $Q_0 = N_1$ equals the tree supremum on both instances \\
G3 & refinement monotonicity, strict on a crafted functional \\
G4 & bridge direction: $Q_0 < 0$ defeats every policy measurably \\
G5 & coupling selection is real: $0$ vs $+1$ at equal marginals \\
G6 & text, retired markers, proof pairing \\
\bottomrule
\end{tabular}
\end{table}

\section{The measure dual of the common-action obstruction}

Let the safe set be $V = \{x: q_{j}(x) \ge 0,\ j = 1,\dots,m\}$ with
$q_{j}$ of class $C^{1}$; let $D$ be a compact-valued upper
semicontinuous disturbance correspondence; let the control set
$U \subseteq \mathbb{R}^{k}$ be compact and \emph{convex}; and let the
dynamics be affine in the control,
$f(x,u,d) = f_{0}(x,d) + G(x)u$, with $f_{0}$ and $G$ continuous. For a
compact information set $B \subseteq V$ (a fibre or estimation tube) with
nonempty active bundle
\[
\mathcal{K}(B) = \{(x,j,d):\ x \in B,\ q_{j}(x) = 0,\ d \in D(x)\},
\]
define the inward drift
$\psi(x,j,d;u) = \nabla q_{j}(x)\cdot[f_{0}(x,d) + G(x)u]$ and the
expected drift $\Phi(u,\mu) = \int_{\mathcal{K}(B)} \psi \,d\mu$ over
Borel probability measures $\mu$ on $\mathcal{K}(B)$.

\begin{theorem}[minimax dual certificate]\label{thm:dual}
Under the standing hypotheses the common safe-action set
$\mathcal{R}^{B}_{V}(B) = \bigcap_{x \in B}\mathcal{R}_{V}(x)$ is empty if
and only if there exist an adversarial measure
$\mu^{*} \in \mathscr{P}(\mathcal{K}(B))$ and a margin $\varepsilon > 0$
with
\begin{equation*}
\max_{u \in U}\,\Phi(u,\mu^{*}) \le -\varepsilon < 0.
\tag{$\star$}
\end{equation*}
Moreover
$\max_{u \in U}\min_{\mathcal{K}(B)}\psi
 = \min_{\mu}\max_{u \in U}\Phi(u,\mu)$.
\end{theorem}
\begin{proof}
$\mathscr{P}(\mathcal{K}(B))$ is weak-$*$ compact and convex (Banach --
Alaoglu), and $\Phi$ is continuous, affine in $u$ (control affinity), and
linear in $\mu$; Sion's minimax theorem (Sion, 1958) applies, and the inner infimum
over measures of a linear functional is attained by the Dirac measures,
so
$\max_{u}\min_{\mu}\Phi = \max_{u}\min_{\mathcal{K}(B)}\psi
 = \min_{\mu}\max_{u}\Phi$.
For the equivalence: $u \in \mathcal{R}^{B}_{V}(B)$ exactly when
$\min_{\mathcal{K}(B)}\psi(x,j,d;u) \ge 0$; the minimum is continuous in
$u$, so nonemptiness of the common safe-action set is equivalent to
$\max_{u \in U}\min_{\mathcal{K}(B)}\psi \ge 0$ (the maximum is attained
on the compact convex $U$), and emptiness is equivalent to the common
value being strictly negative --- which, by the equality, is equivalent to
some $\mu^{*}$ with $\max_{u}\Phi(u,\mu^{*}) < 0$; the margin is the
slack.
\end{proof}

\begin{proposition}[atomic support; $k+1$ witnesses, tight]
\label{prop:sparse}
If $\mathcal{R}^{B}_{V}(B) = \varnothing$, the certificate $\mu^{*}$ may
be chosen atomic with support of at most $k+1$ triples, and the bound is shown tight below in dimensions $k=1$ and $k=2$;
no general-dimensional tightness example is asserted here.
\end{proposition}
\begin{proof}
Write $\psi(x,j,d;u) = v(x,j,d)\cdot u + c(x,j,d)$ with
$v = G(x)^{\top}\nabla q_{j}(x)$ and $c = \nabla q_{j}(x)\cdot f_{0}(x,d)$.
On the compact set $U \times \mathcal{K}(B)$ the function $-\psi$ is
continuous with
$\min_{u}\max_{\mathcal{K}(B)}(-\psi) = 2\eta > 0$ (Theorem~\ref{thm:dual});
a standard finite-cover argument extracts finitely many triples
$(x_{i},j_{i},d_{i})$, $i \le N$, whose drift rows still satisfy
$\min_{u \in U}\max_{i}(-\psi_{i}) \ge \eta$. Min--max duality for the
finite system (linear programming duality over the simplex) yields
$\lambda \ge 0$, $\sum_{i}\lambda_{i} = 1$, with
$\sum_{i}\lambda_{i}\psi_{i}(u) < 0$ on $U$; an extreme such
$\lambda$ is a basic solution of $k+1$ active constraints and is supported
on at most $k+1$ indices (standard theory: Chv\'{a}tal, 1983, chapters 2
and 7; the pairing argument itself is Carath\'eodory's theorem in
$\mathbb{R}^{k}$). The following low-dimensional tightness examples illustrate
the support bound; the plane case can be read through Helly (1923). Tightness: in dimension $k = 2$ the rows
$\{u_{1} \ge 1\}$, $\{u_{2} \ge 1\}$, $\{u_{1}+u_{2} \le 1\}$ on the
compact square $[-10,10]^{2}$ are pairwise intersecting with empty total
intersection (checked on the rational grid by the record), so no two
witnesses suffice; for $k = 1$ the instance of E1 certifies with exactly
two atoms.
\end{proof}

\begin{proposition}[recoveries]\label{prop:recover}
(i) \emph{Polyhedral.} If the drift rows are affine with rational data,
the measure certificate is the $\ell_{1}$-normalized Farkas pair of
Abaee, 2026, An obstruction calculus, Theorem 2: on the two-floor instance
$\{u \le \tfrac25\} \cap \{u \ge \tfrac35\}$ with $U = [0,1]$, the
measure $\lambda = (\tfrac12,\tfrac12)$ makes the expected drift
\emph{constant} $-\tfrac1{10}$, the normalized multiplier pair of the
worked certificate in Section 3.4 there; on the feasible contrast
$\{u \le \tfrac25\} \cap \{u \ge \tfrac15\}$ both sides of
Theorem~\ref{thm:dual} equal $+\tfrac1{10}$ and no certificate exists.
(ii) \emph{Singleton.} For $B = \{x\}$ the dual measures live on the
disturbance set; on the exact convex instance
$\psi(u,d_{1}) = 1-2u$, $\psi(u,d_{2}) = 2u-1$,
$\max_{u}\min_{d}\psi = \min_{\mu}\max_{u}\Phi = 0$, while the
single-worst-disturbance reading $\min_{d}\max_{u}\psi = 1$ is
\emph{not} the dual: the adversarial distribution over disturbances is
load-bearing, and the Isaacs drift is its Dirac-supported inner
infimum.
\end{proposition}

\begin{proposition}[convexity boundary]\label{prop:gap}
The equivalence of Theorem~\ref{thm:dual} fails without control
convexity. With $U = \{0,1\}$ and drift rows
$(-1,+1)$, $(+1,-1)$, the common safe-action set is empty, yet
$\max_{u}\min\psi = -1$ while $\min_{\mu}\max_{u}\Phi = 0$: a strict
minimax gap, and no adversarial measure certifies. The audited finite
action codices of the catalogue notes are non-convex by construction;
there the obstructions are certified combinatorially, and the present
theorem is a companion for the continuous, control-affine layer.
\end{proposition}

\medskip
\noindent\textbf{Shared symbols across the companion papers.} The
companion papers share this vocabulary with paper-local meanings; the
letters below carry meanings here that differ there (including this
paper's own three uses of $K$):
\begin{center}\small
\begin{tabular}{@{}p{0.42\linewidth}p{0.50\linewidth}@{}}
\toprule
\textbf{Here} & \textbf{In the companions} \\
\midrule
$K$: the window/tree depth; $\mathcal{K}(B)$: the contact set; $K_{i}$: the benchmark's carrying capacities & the safe set and kernel $K$ (obstruction calculus); the full-information kernel $K^{T}_{\mathrm{full}}$ (certification); the carrying capacity (forecast ladder) \\
$\mu$, $\mu^{*}$: the adversarial probability measures & Farkas multipliers (obstruction calculus); the certification companion's dual measures \\
$\lambda$: the prior over stored branches; the dual weights & the Farkas multiplier and observer decay rate (obstruction calculus); the certificate weights (certification) \\
$\Gamma_{h}$: the recourse certificate & the certificate margin $\Gamma_{\mathcal{I}}$ (certification); the alpha-vector witness sets $\Gamma^{\Pi}_{k}$ (sufficiency) \\
$\rho$: the interpolation parameter of the DR family & the contamination radius $\rho$ (sufficiency); the mesh relaxation parameter (certification); the net biomass multiplier (applied regime viability) \\
$\beta$: the margin & the stored per-label bound $\beta_{a}$ (certification) \\
$\tau$: the reveal time & the observation time $\tau$ (certification); the delay times (obstruction calculus); the worst-case exit time $\tau^{*}$ (worked systems) \\
$\sigma$: the kernel time variable & the floored exit count $\sigma^{*}$ (worked systems); a noise scale (forecast ladder) \\
\bottomrule
\end{tabular}
\end{center}


\section{The obstruction taxonomy}

The standing axiom adopted from the audit: \emph{a dead state is never an
information failure.} Three classes of nonviability are distinguished and,
on the audited grid, are instantiated side by side (all exact):
\emph{physical} --- the corner $(1,1)$ under two floors is nonviable as a
singleton under full information (no instrument keeps both floors);
\emph{epistemic} --- the audited belief $\{(1,2),(2,1)\}$ has every branch
viable and the joint set empty (failure by merge); \emph{institutional}
--- under the unanimity protocol with agency 2 restricted to vote 1, the
decentralized kernel collapses from $12$ to $0$ (failure by protocol;
full enumeration). The taxonomy disciplines attribution across the
catalogue: mode-blind losses on the same grid are generated entirely at the singleton level, and no count
that mixes the classes is admissible as an epistemic claim.

\section{A calibrated two-stock benchmark}

Leaving the audited grid, the continuous layer needs a worked system with
one fixed control set throughout; the two-stock system below is the
running instance of the continuous layer --- the caps-fibre benchmark
shared with the worked-systems companion. Two renewable stocks
($r_{1} = 1$, $r_{2} = \tfrac45$, $K_{1} = K_{2} = 10$, cross-competition
$\tfrac1{20}$, shocks $|d_{i}| \le \tfrac1{10}$), floors
$x_{i} \ge 2$, and the single control set
$U = \{u \ge 0:\ u_{1} + u_{2} \ge 2\}$ (a demand floor; no capacity
constraint is imposed). On the aggregate-yield fibre
$Y = x_{1} + x_{2}$ within the floors, the active-floor harvest caps are
\[
\mathrm{cap}_{1}(Y) = \tfrac32 - \tfrac{Y-2}{10},\qquad
\mathrm{cap}_{2}(Y) = \tfrac{59}{50} - \tfrac{Y-2}{10},
\]
so the fibre is viable exactly when
$\mathrm{cap}_{1}(Y) + \mathrm{cap}_{2}(Y) \ge 2$, i.e.
\begin{theorem}[critical aggregate]\label{thm:benchmark}
$Y^{*} = \tfrac{27}{5}$: aggregate fibres with $Y \le \tfrac{27}{5}$ are
viable and those with $Y > \tfrac{27}{5}$ are obstructed, the value being
exact ($\mathrm{cap}_{1} + \mathrm{cap}_{2} = 2$ at $Y^{*}$). At
$Y = 5$ the fibre is viable with witness $u = (\tfrac65,\tfrac45)$; at
$Y = 6$ it is obstructed (cap sum $\tfrac{47}{25} < 2$) with dual
certificate $\mu^{*} = (\tfrac12,\tfrac12)$ on the two active-floor
atoms --- two points, within the $k+1 = 3$ bound --- and margin
$\varepsilon = \tfrac{3}{50}$. The fibre criterion itself --- viability
exactly on the cap-sum inequality, audited per aggregate with radius
bounds --- is carried in the worked-systems companion (Abaee, 2026, Worked systems for the obstruction calculus);
this section contributes the two-stock derivation and the fibre-width
formula.
\end{theorem}
The reading is the opposite of the intuition that abundance certifies
safety: cross-competition tightens the harvest caps as the aggregate
grows, so richer fibres obstruct first --- aggregate yield is a failing
indicator in both directions. The fibre-width formula makes the
aggregation rule exact: the $Y$-fibre within the floors is
$x_{1} \in [2,\,Y-2]$, of width $\Delta x_{1}(Y) = Y - 4$, so a linear
index cannot separate states closer together than the floor buffer ---
the composite index is certifying exactly when the width stays inside the
buffer, and the certainly-safe reporting of the applications companion
(Abaee, 2026, Robust viability of the 2J3KL limit reference point) governs otherwise.

\section{The dynamic envelope}

A natural candidate formulation of a dynamic certificate suggests itself: an adversarial measure curve $t \mapsto \mu^{*}_{t}$ whose expected-margin envelope $Q(t)$ decreases at a certified rate, governed by an information-set operator acting on the state. The candidate fails, for three reasons that shape the formulation developed below; this section records the reasons, formulates the envelope, and proves what exact methods reach. The residue remains conjectural.

\subsection{Three obstructions to the candidate}

\emph{(i) A Bellman object posed as a pointwise Hamiltonian.} The candidate operator maps a
value along trajectories to an expectation over an information set of
measures --- a Bellman-type, hence nonlocal, object --- while dressed as a
pointwise differential operator. Restricted-information games do not
produce pointwise Hamiltonians in the state; their well-posed home is the
information state (here, the belief), on which the object is a dynamic
program. The nonlocality is not itself the obstruction; the pointwise differential form is.

\emph{(ii) An inconsistent measure curve.} A curve $t \mapsto
\mu^{*}_{t}$ of measures carries no requirement that $\mu^{*}_{t}$ be the
marginal at time $t$ of one measure on disturbance paths. Without that
consistency the envelope has no time consistency: nothing connects $Q(t)$
to $Q(t')$ for $t' > t$, so no backward recursion and no comparison can
even be posed. The formulation developed below replaces the curve by a single adversarial prior $\mu^{*}$ on disturbance paths; conditional expectation then supplies
the martingale structure without any PDE.

\emph{(iii) Expectation versus trajectory.} For the certificate direction
the selection argument is not needed at all: an expectation of a bounded
functional that is strictly negative exhibits a positive-measure set of
realizations on which the functional is strictly negative. Selection
theorems are required only for the converse value-attainment direction,
which the certificate does not use.

\subsection{The envelope formulation}

\begin{lemma}[expectation-to-realization bridge]\label{lem:bridge}
Let $\mu^{*}$ be a probability measure on disturbance paths compatible with
the information structure, let $\pi$ be nonanticipative, and let the
safety functional $F^{\pi}(\omega)$ --- the realized weighted facet margin
at the horizon --- be bounded below and $\mathcal{F}_{T}$-measurable. If
$\mathbb{E}^{\mu^{*}}[F^{\pi}] < 0$, then
$\mu^{*}(\{\omega: F^{\pi}(\omega) < 0\}) > 0$: some compatible disturbance
realization violates the margin, and the set of such realizations carries
positive $\mu^{*}$-measure.
\end{lemma}
\begin{proof}
Nonanticipativity of $\pi$ makes $F^{\pi}$ an $\mathcal{F}_{T}$-measurable
functional of the disturbance path, and boundedness below makes its
expectation well-defined. If $\mu^{*}(F^{\pi} < 0) = 0$, then
$F^{\pi} \ge 0$ holds $\mu^{*}$-almost surely, so
$\mathbb{E}^{\mu^{*}}[F^{\pi}] \ge 0$ --- contradicting the hypothesis. The
realizations of the positive-measure set are compatible by construction,
since $\mu^{*}$ is supported on the compatible set.
\end{proof}

\begin{definition}[adversarial-prior envelope]\label{def:envelope}
Fix a horizon $T$, a bounded safety functional $F$ (positive values safe),
and one consistent adversarial prior $\mu^{*}$ on disturbance paths. The
envelope is
\[
Q(t) \;=\; \operatorname*{ess\,sup}_{\pi} \;
\mathbb{E}^{\mu^{*}}\bigl[\, F^{\pi} \;\big|\; \mathcal{F}_{t} \,\bigr],
\]
the essential supremum over nonanticipative policies of the conditional
expected safety functional. For each fixed $\pi$ the process
$t \mapsto \mathbb{E}^{\mu^{*}}[F^{\pi} \mid \mathcal{F}_{t}]$ is a
martingale; the envelope is its essential supremum over policies, and its
pointwise optimization over \emph{complete} policies can select
different past actions on different later cells. It need not obey a
Bellman recursion unless committed action history is included in the
state and only future actions are optimized.
\end{definition}

\begin{theorem}[finite well-posedness, comparison, exactness]
\label{thm:envelope-finite}
In the finite setting --- disturbance set $\Omega$ finite, information
states $b$ drawn from a finite partition together with their
committed action prefixes (or a proved Markov-sufficient state encoding
them), admissible actions finite, data rational --- define the
one-step continuation operator
\[
\Phi V(b) \;=\; \max_{a}\; \mathbb{E}^{\mu^{*}}\bigl[\, V(b') \,\big|\, b, a \,\bigr],
\]
where $b'$ records the posterior and the committed prefix
extended by $a$ after observation $y$; the maximum concerns only
not-yet-chosen actions. Bare posterior cells alone need not be a
sufficient state for path-dependent terminal $F$. Then: (i) $\Phi$ is order-preserving,
$V \le W \Rightarrow \Phi V \le \Phi W$ pointwise; (ii)
with committed prefixes the initial attainable value $Q(0)$ is
computed exactly by the finite backward induction \(\Phi^N g\),
in rational arithmetic; this does not identify intermediate values
with the pointwise essential supremum over complete policies in
Definition~\ref{def:envelope}; and (iii) $Q(0) < 0$ certifies that for every admissible policy
the set of disturbance realizations defeating it carries positive
$\mu^{*}$-measure.
\end{theorem}
\begin{proof}
(i) If $V \le W$ pointwise, then for each $(b, a)$ the conditional
expectation of $V(b')$ is at most that of $W(b')$ (monotonicity of
expectation over the same kernel), and the maximum over the same action set
preserves the inequality. (ii) By induction on the remaining horizon:
$\Phi^{0} g = g$ is exact by definition; the induction step applies $\Phi$
once, and each evaluation is a finite weighted sum of previously computed
values, so rational data propagate exactly through the recursion.
Well-posedness of the committed-prefix recursion follows from
the same induction. For a counterexample to posterior-only recursion,
choose a blind first action $a_1\in\{0,1\}$ and then reveal a
fair bit $s\in\{0,1\}$, with terminal reward
$G(s,a_1)=\mathbf{1}[a_1=s]$. Every feasible complete policy has
initial value $1/2$; at the revealed cell the pointwise essential
supremum over \emph{complete} policies is $1$ for each bit, by
retrospectively choosing two different first actions. A Bellman
continuation retaining the first action returns $1/2$, as required. (iii) $Q(0) < 0$ bounds every policy's
conditional expectation: for the optimal policy the expectation is
$Q(0) < 0$, and for every suboptimal $\pi$ the expectation is smaller
still, so Lemma~\ref{lem:bridge} applies to each $\pi$ and yields the
positive-measure defeating set.
\end{proof}

\begin{proposition}[identification with the recourse certificate]
\label{prop:recourse-id}
One step deep, with stored facet margins $\beta_{\ell}$, control kernels
$k_{\ell}(\sigma)$, prior $\lambda$ over the stored branches, and control
set $U$, the envelope evaluation of
Definition~\ref{def:envelope} --- the policy's optimal ceiling on the
expected weighted facet sum --- is exactly the support-function certificate
\[
\Gamma_{h}(\lambda) \;=\; \sum_{\ell} \lambda_{\ell} \beta_{\ell}
\;+\; \int_{0}^{\tau} h_{U}\Big(\sum_{\ell} \lambda_{\ell} k_{\ell}\Big)
+ \sum_{\ell} \int_{\tau}^{T} h_{U}\big(\lambda_{\ell} k_{\ell}\big),
\]
with $h_{U}(g) = \max_{v \in U} g^{\top} v$ the support function of $U$,
of the obstruction calculus (its Section 3.7; complete proof in its
Supplementary Material, S1). The recourse certificate is the dynamic
envelope's one-step evaluation.
\end{proposition}
\begin{proof}
Variation of constants expands each realization of the weighted facet sum
into the stored $\beta$-term plus the window pairing (one common signal on
$[0, \tau)$) plus the branch pairings (separate post-reveal signals). The
supremum of a linear functional over the signal sets factorizes across the
time intervals, and each interval's supremum of a realized pairing
$g^{\top} v$ over $v \in U$ is the support value $h_{U}(g)$; collecting
terms gives exactly $\Gamma_{h}(\lambda)$ as the ceiling of every policy's
realized weighted facet sum. Negativity of the ceiling certifies by
Lemma~\ref{lem:bridge}: the realized weighted facet sum of every policy is
at most $\Gamma_{h}(\lambda) < 0$, so some facet is violated on a
positive-measure set of realizations.
\end{proof}

\begin{corollary}[the three-branch instance, certified]\label{cor:three-branch}
On the braking instance of the obstruction calculus (three branches,
weights $\lambda = (\tfrac{3}{8}, \tfrac{5}{16}, \tfrac{5}{16})$, stored
margin $\beta = \tfrac{16}{25} - \tau - 1$): the pooled window kernel
vanishes ($\sum_{\ell} \lambda_{\ell} n_{\ell} = 0$), each branch's
post-reveal pairing contributes its support value
$\lambda_{\ell} \int_{\tau}^{\tau+1} (\tau + 1 - \sigma)\, d\sigma =
\lambda_{\ell}/2$, and the envelope evaluation is
$\Gamma_{h}(\tau) = \tfrac{7}{50} - \tau$: negative exactly for
$\tau > \tfrac{7}{50} = 0.14$, where the bridge converts it into the
positive-measure defeating set for every blind-window policy;
hold-then-brake is safe exactly to $\tau = \tfrac{7}{50}$, so the
certificate is tight. The verification script evaluates the identity
exactly at $\tau \in \{\tfrac{1}{10}, \tfrac{7}{50}, \tfrac{1}{5}\}$.
\end{corollary}

\begin{conjecture}[the open residue in continuous time]\label{conj:envelope}
The continuous-time residue of the envelope formulation: (i) an evolution
law for the adversarial prior along the review sequence --- a consistent
coupling of the one-step priors, in the spirit of martingale-transport
couplings (Beiglb\"ock, Henry-Labord\`ere, and Penkner, 2013) --- under which the envelope obeys a dynamic-programming equation
in continuous time; (ii) the limit of the information-state recursion under
partition refinement, which is where a comparison principle must live if a
differential statement is to be recovered at all; and (iii) value
attainment (the selection direction) in the general filtered setting. The
static identification of Proposition~\ref{prop:recourse-id} is exact and is
not part of the conjecture.
\end{conjecture}

\section{The review-sequence envelope: exact couplings}

The envelope of the preceding section is one review deep. This section certifies
the two properties its continuous-time residue must preserve --- coupling
consistency and refinement monotonicity --- on exact review-sequence
instances; everything is rational arithmetic, regenerated by the script.

\subsection{The instance}

Two windows, one hidden branch. The branch is $s \in \{+1, -1\}$, prior
$\tfrac{1}{2}$ each; the state starts at $z_0 = 2$. At each review the
policy acts $u_1, u_2 \in \{-2, 0, 2\}$; between reviews the adversarial
disturbance acts $d_1, d_2 \in \{-1, 1\}$ (one consistent prior: the
product measure, $\tfrac{1}{2}$ per value per window). The state obeys
$z_1 = z_0 + s\, u_1 + d_1$ and $z_2 = z_1 + s\, u_2 + d_2$; the review
after window 1 observes $z_1$ exactly. The safety functional is the
expected two-sided margin
\[
F \;=\; \mathbb{E}_{s}\bigl[\, 2 - |z_{2}(s)| \,\bigr],
\qquad z_{2}(s) = z_0 + s(u_1 + u_2) + d_1 + d_2,
\]
bounded on the finite instance, so the bridge lemma applies verbatim.

\begin{proposition}[consistent coupling and the tower identity]
\label{prop:tower}
The one-step adversarial priors couple to the product measure on
$(d_1, d_2)$, and for \emph{every} nonanticipative policy on the instance
--- all $3^{4} = 81$ policy trees over the reachable $z_1$ cells --- the
path-enumerated expectation and the iterated-conditional (tower)
evaluation of $\mathbb{E}[F]$ agree exactly. Each fixed policy's
conditional-expectation process along the review filtration is therefore
a martingale, and the envelope
$Q_0 = \operatorname*{ess\,sup}_{\pi} \mathbb{E}[F \mid \mathcal{F}_0]$
is computed by backward induction: $Q_0 = \Phi^{2} g$ equals the
policy-tree supremum ($= \tfrac{1}{2}$, at $u_1 = \pm 2$).
\end{proposition}
\begin{proof}
Consistency is by construction: the product measure has the one-step
priors as its marginals. For the tower identity, condition on
$\mathcal{F}_1$ (the review): within each $z_1$ cell the inner
conditional expectation is a finite weighted average, and the outer
average over the cells reconstructs the path sum because the cells
partition the paths --- an induction on the two-node filtration, exact
in rationals. The martingale property is the tower identity restated
along the filtration. Backward induction computes
$\Phi^{2} g$: $\Phi g$ at each cell is a maximum of three rational
expectations, and $\Phi$ is applied once more over $u_1$; the script
checks the result against the enumeration of all $81$ trees.
\end{proof}

\begin{proposition}[information design and recourse strictly pay]
\label{prop:strict}
On the instance: the envelope value is $Q_0 = \tfrac{1}{2}$, attained by
$u_1 = -2$ (any $u_1 = \pm 2$ by symmetry) with the recourse
$u_2(z_1) = (0, -2, +2, +2)$ at $z_1 = (-1, 1, 3, 5)$; every open-loop
policy and the do-nothing policy attain only $0$. The strict gap
$\tfrac{1}{2} > 0$ is the certified content of the correction: the first
review's action \emph{designs the information} ($u_1 = -2$ makes $z_1$
separate the branch exactly), and the second \emph{exploits it}
(brake toward the origin on the revealed branch); at $z_1 = 5$ the
revealed state is unrecoverable and the cell value is $-1$.
\end{proposition}
\begin{proof}
The four singleton cells at $u_1 = -2$ carry values
$Q_1 = (1, 1, 1, -1)$ at weights $\tfrac{1}{4}$ each: $Q_0 = \tfrac12$.
The open-loop maximum is the better of nine fixed pairs, computed
exactly as $0$; the do-nothing pair $(0,0)$ also evaluates to $0$.
The cell values and the maximizing $u_2$ are direct rational
computations, reproduced by the script.
\end{proof}

\begin{proposition}[refinement is monotone]\label{prop:refine}
On the three-state matching instance --- hidden $x \in \{1,2,3\}$, prior
$\tfrac13$ each, margin $-|x - a|$, action $a$ per observation block ---
the envelope is non-decreasing under partition refinement and strictly
increasing on the instance's final step:
$Q(\{\{1,2\},\{3\}\}) = Q(\{\{1\},\{2,3\}\}) = -\tfrac13
< 0 = Q(\{\{1\},\{2\},\{3\}\})$.
\end{proposition}
\begin{proof}
Every policy measurable under a coarser partition is measurable under a
refinement, so the supima are ordered. Strictness: under the fine
partition the policy matches each state exactly ($a = x$, value $0$);
under either coarse partition two states share an action and the best
shared action leaves one mismatch, value $-\tfrac13$. This monotone,
bounded family is the direction the partition-refinement limit of the
residue conjecture must take.
\end{proof}

The residue conjecture now reads sharper: the continuous-time coupling
law must reduce, on review sequences with finite disturbance sets, to
the product coupling certified here; and the information-state recursion
must converge along refinements with the monotonicity certified here.

\section{The general finite-sequence envelope}

The instance of the previous section is the $K = 2$ case of a general
statement. This section states the general theorem over finite review
sequences and proves it in full; the two-window instance is the $K = 2$ case, and the general argument uses only finite combinatorics and the linearity of expectation.

\subsection{Setup}

Fix a number of windows $K \ge 1$. Each window $k$ carries a finite
disturbance set $D_{k}$; the disturbance paths form
$\Omega = D_{1} \times \cdots \times D_{K}$, a finite set. A
\emph{consistent adversarial prior} is a product measure
$\mu^{*} = \mu^{*}_{1} \otimes \cdots \otimes \mu^{*}_{K}$ on $\Omega$;
consistency --- each marginal $\mu^{*}_{k}$ is the time-$k$ law of the one measure $\mu^{*}$ --- is exactly the requirement whose absence makes a backward recursion ill-posed. Reviews divide the sequence: the review
filtration is an increasing sequence of partitions
$\mathcal{P}_{1} \preceq \mathcal{P}_{2} \preceq \cdots \preceq
\mathcal{P}_{K}$ of $\Omega$, each $\mathcal{P}_{k}$ consisting of the
cells of the information available at review $k$ (so
$\mathcal{P}_{k}$ refines into $\mathcal{P}_{k+1}$), with
$\mathcal{P}_{K}$ the singleton partition; $\mathcal{P}_{1}$ may be the
trivial partition (a blind first window). Window $k$ has a finite action
set $A_{k}$; a \emph{nonanticipative policy} is
$\pi = (\pi_{1}, \dots, \pi_{K})$ with $\pi_{k}$ constant on the cells
of $\mathcal{P}_{k}$ --- the filtration carries the whole observation
history, so cell-constancy is adaptedness. The data are completed by a
bounded safety functional $G: \Omega \times A_{1} \times \cdots \times
A_{K} \to \mathbb{R}$ (positive values safe, per the sign convention of
this paper), rational throughout; the realized functional of $\pi$ is
$G^{\pi}(\omega) = G(\omega, \pi_{1}[\omega]_{1}, \dots,
\pi_{K}[\omega]_{K})$, where $[\omega]_{k}$ is the
$\mathcal{P}_{k}$-cell of $\omega$.

\begin{theorem}[general finite-sequence envelope]\label{thm:general}
Under the setup above, use the recursion only on cells
$C$ with $\mu^*(C)>0$; omit null children from its sums and leave
null-node conditional values undefined (or assign arbitrary values
that never affect the root). In particular the initial node is
positive-mass. Then: \emph{(i) tower and martingale.} For every
nonanticipative policy $\pi$,
\[
\mathbb{E}^{\mu^{*}}[G^{\pi}] \;=\;
\mathbb{E}^{\mu^{*}}\bigl[\, \mathbb{E}^{\mu^{*}}\bigl[\, \cdots
\mathbb{E}^{\mu^{*}}[\, G^{\pi} \mid \mathcal{P}_{K} \bigr] \cdots
\mid \mathcal{P}_{2} \bigr] \mid \mathcal{P}_{1} \bigr],
\]
and $M^{\pi}_{k} = \mathbb{E}^{\mu^{*}}[G^{\pi} \mid
\mathcal{P}_{k}]$ is a $\mu^{*}$-martingale along the filtration.
\emph{(ii) Dynamic programming.} Define the node values $N$ on the
review tree --- nodes are pairs (cell, action prefix) --- by
$N_{K}\bigl(\{\omega\},\, a_{1:K-1}\bigr) = \max_{a \in A_{K}}
G(\omega, a_{1:K-1}, a)$ and, for $k < K$, the sum running over the
children $C' \subseteq C$ of $\mathcal{P}_{k+1}$ over positive-mass children only, with the child weights
$w(C') = \mu^{*}(C')/\mu^{*}(C)$,
\[
N_{k}\bigl(C,\, a_{1:k-1}\bigr)
\;=\; \max_{a \in A_{k}} \;
\sum_{C'} w(C')\, N_{k+1}\bigl(C',\, a_{1:k-1}a\bigr),
\]
Then the envelope value $Q_{0} := \sup_{\pi}
\mathbb{E}^{\mu^{*}}[G^{\pi}]$ equals $N_{1}$ evaluated on the
$\mathcal{P}_{1}$-cells, averaged by their prior weights
$Q_{0} = \sum_{C \in \mathcal{P}_{1}:\mu^*(C)>0} \mu^{*}(C)\,
N_{1}(C)$ (a single node if $\mathcal{P}_{1}$ is trivial); the supremum
is attained by a nonanticipative policy read off the maximizing
arguments of $N$. \emph{(iii) Refinement monotonicity.} If a second
filtration $\mathcal{P}'$ refines $\mathcal{P}$ stage-wise, then the
induced envelope values satisfy $Q_{0}(\mathcal{P}) \le
Q_{0}(\mathcal{P}')$. \emph{(iv) Exactness and the obstruction
direction.} All node values are rational; if $Q_{0} < 0$, then for
\emph{every} nonanticipative $\pi$ the set
$\{\omega: G^{\pi}(\omega) < 0\}$ carries positive $\mu^{*}$-measure.
\end{theorem}
\begin{proof}
\emph{(i)} Each $G^{\pi}$ is $\mathcal{P}_{K}$-measurable (a finite
function of the path and the cell-constant actions), and repeated
conditioning is an identity on finite probability spaces; the tower
formula is that identity read backwards, and the martingale property is
$M^{\pi}_{k} = \mathbb{E}[M^{\pi}_{k+1} \mid \mathcal{P}_{k}]$, which
is the defining property of conditional expectation applied twice.
\emph{(ii)} Positive-mass conditioning makes each weight
well-defined and null children contribute zero to the root. Proceed
by backward induction on the stage. For $k = K$ the node
value is the best terminal action, and $G^{\pi}$ at a singleton cell is
$G(\omega, \cdot)$ evaluated at the realized prefix and $\pi_{K}$, so
$\max_{a_{K}} G = N_{K}$ pointwise. For the step, fix a
$\mathcal{P}_{k}$-cell $C$ and an action prefix, and split the policies
by their stage-$k$ action: every policy compatible with the prefix is a
choice $a \in A_{k}$ together with a continuation policy for the
conditioned sub-instance on $\{C\}$-successor cells, whose filtration
$\bigl(\mathcal{P}_{j}\bigr)_{j > k}$ restricted to the successors
refines to singletons at $K$. By (i) applied at the conditioning,
\[
\mathbb{E}^{\mu^{*}}[G^{\pi} \mid C] \;=\;
\sum_{C' \subseteq C} \frac{\mu^{*}(C')}{\mu^{*}(C)}\;
\mathbb{E}^{\mu^{*}}\bigl[G^{\pi} \mid C'\bigr],
\]
the conditional law of the successor cell given $C$ is the ratio of
prior masses (exogeneity: the prefix does not move $\mu^{*}$), and the
induction hypothesis --- applied to each conditioned sub-instance,
which has one fewer window --- identifies the supremum of the inner
conditional expectations over continuations of $(a_{1:k-1}, a)$ with
$N_{k+1}(C', a_{1:k-1}a)$. Taking the outer maximum over $a$ gives
$N_{k}$; summing over the $\mathcal{P}_{1}$-cells with their prior
weights and applying (i) once more gives $Q_{0} = \sum_{C}
\mu^{*}(C) N_{1}(C)$. Attainment: each $N$ node is a maximum over a
finite set; selecting a maximizing argument at every node and pasting
along the tree yields a nonanticipative policy --- cell-constancy holds
because the selection is made per node --- and its expectation is
exactly the recursion's value. \emph{(iii)} If $\mathcal{P}'$ refines
$\mathcal{P}$ stage-wise, every $\mathcal{P}$-cell-constant policy is
$\mathcal{P}'$-cell-constant, so the $\mathcal{P}'$-supremum ranges
over a superset of policies and cannot be smaller; the inequality for
the full sequences follows stage-wise from the one-step inclusion.
\emph{(iv)} The recursion's only operations are maxima and
$\mu^{*}$-weighted averages of rationals, so every $N$ is rational. If
$Q_{0} < 0$, then by (ii) every policy satisfies
$\mathbb{E}^{\mu^{*}}[G^{\pi}] \le Q_{0} < 0$, and the
expectation-to-realization bridge (Lemma~\ref{lem:bridge}) converts strict negative expectation into a
positive-measure set of strictly negative realizations.
\end{proof}

\begin{remark}[the recourse certificate and the residue]\label{rem:after-general}
Three consequences. \emph{First}, the identification with the recourse
certificate is the one-step case of Theorem~\ref{thm:general}: with
stored facet margins the node recursion collapses to the
support-function evaluation $\Gamma_{h}$ (the identification
proposition), so the calculus's recourse
certificate is the envelope's $K$-node tree read at depth one.
\emph{Second}, the one-step priors $\mu^{*}_{k}$ do not determine the
coupling $\mu^{*}$: two consistent couplings with identical one-step
marginals can give different node values (the script certifies a
two-window example where the product coupling yields $0$ and the
comonotone coupling yields $+1$ on the same functional). The residue
conjecture's ``evolution law for the adversarial prior'' is therefore
precisely a \emph{coupling-selection} problem --- choose the consistent
coupling as a function of the state and review history --- which is the finite-space form of the martingale-transport structure recorded in the conjecture. \emph{Third}, refinement monotonicity (iii) is the
finite-space grounding of the partition-limit clause: the instance
family of the previous section, monotone and bounded, now carries a
general ordering, and the continuous-time residue is exactly the
question of which couplings and which refinement limits survive the
passage to continuous disturbance sets. The distributionally robust
interpolation of the probabilistic companion (Abaee, 2026, Probabilistic sufficiency) holds the
adversary to a mixture family: on the two-window instance its value is
exactly $\\rho$ at parameter $\\rho$ --- the interpolation contains the
coupling extremum at $\\rho = 1$, preserves the one-step marginals at
every $\\rho$, and leaves the coupling-selection difficulty orthogonal to
the mixture parameter (joint check archived).
\end{remark}

\section{Mechanized verification}

The Lean~4 module \texttt{Formalizations/Minimax\_Dual} proves
conditional soundness of a supplied finite measure certificate and
specific algebraic examples, not the Sion converse, sparse bound or
dynamic envelope. A current comment-stripped count including Unicode names finds \(22\)
\texttt{theorem} declarations and no \texttt{lemma} declarations.
The historical report of \(21\) is an undercount: an ASCII-only name
filter reproduces it by omitting \texttt{\ensuremath{\psi}single\_sum}.
The original historical counting command has not been recovered; this
is a reproducible explanation, not proof of its exact implementation.
The declared-pin default-target build now passes in GitHub Actions. The module imports only the
project prelude. The module builds as part of the project's \(60\)-job build, all of
which pass.

\noindent\textbf{No admitted gaps.} In Lean an unfinished proof appears as the axiom
\texttt{sorryAx} in the axiom footprint of every declaration depending on it. That
footprint was inspected: \texttt{sorryAx} is absent, and a source scan with comments
stripped finds no \texttt{sorry}, \texttt{admit} or \texttt{axiom} declaration anywhere in
the formalization layer. The named Minimax declarations checked here use only those standard classical axioms or none; the wider project also has the P3 native-evaluation dependencies disclosed below.
\noindent\textbf{Current Lean provenance.} The committed project pin \texttt{v4.34.1} was rebuilt in GitHub Actions (default Lake target: \(60/60\) jobs, exit~0; \url{https://github.com/MIKEAA2020/general-sustainability/actions/runs/36946730855}). The exact project source and two footprint-check \texttt{.lean} files are archived with SHA-256 manifests. A comment-aware scan of all \(62\) source files finds no \texttt{sorry}, \texttt{admit}, explicit \texttt{axiom} or \texttt{constant} declaration. Twenty-one named footprints contain no \texttt{sorryAx}; the aggregate-import environment exposes three source-level \texttt{native\_decide} uses in the P3 support/value examples, yielding nine distinct generated axiom names under this pin. These are executable computational certificates that extend kernel-only trust to the Lean compiler/runtime; they are not admitted gaps and must not be described as zero axioms or as kernel-only proofs. Three preserved historical modules outside the aggregate import closure are not certified by the default build. The pinned declared-exceptions gate passed in GitHub Actions (run 36990482599): only the nine exact generated names from three approved P3 source sites are permitted; any change to those sites or names requires separate review and re-approval. The separate strict zero-exception check exits 2 by design, while the build itself passed 60/60; the two verdicts must not be conflated. The v4.14.0 comparison of the \emph{current} source fails in the prelude and does not reproduce the old source run.


\noindent\textbf{Scope.} What is machine-checked is the dual-certificate core stated over
an ordered-field interface, so it holds in any model of that interface. Numerical
instantiations, and any continuous-time statements, are not covered.

\section{Related work}
\label{priorart}

The question this paper answers --- is there a single instrument that is safe against every
branch of an information set, and if not, what certifies that there is not --- has a long
history under other names. This section fixes the position against each of the four
neighbouring programmes. The short version: the \emph{question} is not new, and one of these
programmes already answers it in a different form; what is new is the \emph{form of the
answer} (a finitely supported measure with a tight support bound) and the \emph{precise
location of the boundary} at which that form stops existing.

\subsection{The honest baseline: this is the discriminating-kernel question}

The viability-theoretic treatment of two-player, state-constrained games is built on the
\emph{discriminating kernel} (Aubin, 1991; Aubin and Catt\'e, 2002; Cardaliaguet, Quincampoix
and Saint-Pierre, 1994, 2007): the largest closed subset of the constraint set that is a
discriminating domain, characterised as Victor's victory domain and, in the discrete and
continuous settings alike, computed by a descending-kernel algorithm. The discriminating
kernel is the two-player analogue of the viability kernel, and its algebraic characterisation
is as the largest, smallest, or \emph{unique minimax} --- bilateral --- fixed point of an
adequate map on pairs of subsets (Aubin and Catt\'e, 2002).

\textbf{The static local obstruction is not a global-kernel test.}
Theorem~\ref{thm:dual} certifies emptiness of the common
\emph{instantaneous} safe-action set at the declared information
state under its compact-control hypotheses. A global discriminating
or viability kernel additionally requires a horizon/invariance
argument. For example, on $[0,1]$ with $\dot x=1$ and $U=\{0\}$,
$x=0$ has an inward-pointing instantaneous action, but every
trajectory eventually exits and the infinite-horizon kernel is
empty. The established global-kernel programme is therefore related
prior art, not a logically equivalent restatement of this static dual.

Four things are not in that programme, and they are what this paper claims.

\begin{enumerate}
\item \textbf{The certificate is a measure, not a set.} The discriminating-kernel programme
  returns a \emph{subset of state space}, characterised by a fixed-point property and computed
  by a set iteration. Theorem~\ref{thm:dual} returns a \emph{probability measure} on the
  active boundary--disturbance bundle, finitely supported, whose existence is equivalent to
  emptiness of the \emph{local common safe-action set} under its hypotheses.
  This does not by itself decide membership in a global invariant kernel.
\item \textbf{The support bound, and its tightness.} Proposition~\ref{prop:sparse} gives a
  certificate supported on at most $k+1$ points, where $k$ is the \emph{control} dimension,
  with tight examples shown here for $k=1,2$. Note the scaling: the witness count is governed by the dimension
  of the instrument space, not of the state space. Nothing in the fixed-point characterisation
  of kernels yields a finite witness of this kind.
\item \textbf{The convexity boundary is located exactly where the kernel programme puts it,
  from the other side.} Cardaliaguet, Quincampoix and Saint-Pierre (1994, Proposition~2.3)
  show that under suitable convexity assumptions the discriminating kernel is
  \emph{convex}. Proposition~\ref{prop:gap} here is the complementary negative: without
  convexity the duality fails outright, and an explicit two-action instance exhibits a strict
  minimax gap in which \emph{no} measure certifies. The two results are consistent and
  together say that the convexity hypothesis is load-bearing in the strongest possible sense
  --- it is simultaneously what makes kernels convex and what makes measure duality sound.
\item \textbf{The certificate is one-sided and observation-constrained, by design.} The
  discriminating-kernel programme computes a victory domain in both directions. This paper
  certifies only the negative direction (no common safe action exists) for an information set
  under partial observation. That restriction is what buys items~1 and~2; it is a narrowing,
  not an improvement, and is stated as such.
\end{enumerate}

\subsection{Farkas, and linear-programming duality}

The polyhedral case of Theorem~\ref{thm:dual} recovers the $\ell_{1}$-normalised Farkas
multiplier pair exactly (Farkas, 1902), and the classical alternative is the mechanism on
which the whole certificate rests. Farkas' lemma is an alternative for a finite system of
linear inequalities: exactly one of a primal feasibility and a dual infeasibility certificate
holds.

The step beyond it is the measure-valued, non-polyhedral form. The classical alternative is
finite-dimensional and static; here the dual object ranges over probability measures on a
boundary--disturbance bundle, the emptiness being certified is that of a set of
\emph{simultaneously} safe actions across an information set rather than of a polyhedron, and
the recovery of the Farkas pair in the polyhedral case
(Proposition~\ref{prop:recover}) is a consistency check on the generalisation, not the
generalisation itself.

\subsection{Isaacs, and the value-function programme}

The static picture is completed by the Isaacs drift condition (Isaacs, 1965). Where the
classical Hamilton--Jacobi--Isaacs theory produces a \emph{value} as a function of the state
under full observation, the object here is a one-directional refutation --- a proof that no
common instrument exists --- computed over an information set. The singleton case of
Theorem~\ref{thm:dual} recovers the Isaacs drift with the adversarial \emph{distribution}
over disturbances load-bearing (Proposition~\ref{prop:recover}); the distribution is what
carries the argument, and that is why the singleton case is not merely the classical
condition restated.

\subsection{Differential games under incomplete information, on the space of measures}

The closest single neighbour is the programme that lifts a zero-sum differential game with
\emph{imperfect} information about the state onto the space of probability measures.
Cardaliaguet and Quincampoix (2008) study the case in which the players know only a
distribution on the initial state, and establish existence of the value through a uniqueness
result for Hamilton--Jacobi--Isaacs equations stated on the space of measures. Both that
work and this one are measure-valued formulations of a game under incomplete information.

The aims differ. That programme establishes that a \emph{value exists}, by analytic means, in
an infinite-dimensional setting. This paper produces a \emph{finite certificate}: a
finitely supported measure on $k+1$ points, checkable in exact arithmetic, whose existence
refutes the common-action property. An existence-and-uniqueness result for a value does not
yield a finite witness, and the sparsity and tightness of
Proposition~\ref{prop:sparse} have no analogue there. Conversely, nothing here establishes
the existence of a value, and the paper does not claim to.

\subsection{Minimax equality}

That minimax equality holds for quasi-concave--convex payoffs with appropriate compactness
and semicontinuity is classical (Sion, 1958), generalising von Neumann. The equality of
Theorem~\ref{thm:dual} in the convex case can be read as an instance of that classical
circle of results, and this paper makes no claim to have originated the equality.

What the classical theorems do not supply is the sparse witness. Sion's theorem yields
\emph{existence}; it says nothing about the support size of a certifying measure, and the
$k+1$ bound of Proposition~\ref{prop:sparse} --- together with the construction showing it is
saturated --- is not recoverable from it. Symmetrically, Proposition~\ref{prop:gap} is
consistent with the classical theory: the minimax equality is precisely what
\emph{fails} once the convexity hypothesis is dropped.

\subsection{Verified numerics, and what is claimed}

The certificate is checked in exact rational arithmetic, with independent primal and dual
witnesses; the solver output is a confirmation, not a proof object. This is a statement about
the pipeline, not a claim over validated-integration methods, which address problems this
paper does not.

\medskip

\noindent \textbf{Claimed:} the measure form of the obstruction and its tight $k+1$ support
bound (Theorem~\ref{thm:dual}, Proposition~\ref{prop:sparse}); the two recoveries, Farkas and
Isaacs (Proposition~\ref{prop:recover}); the exact location of the convexity boundary with an
explicit gap instance (Proposition~\ref{prop:gap}); the dynamic, review-sequence and
finite-sequence envelope extensions with the tower identity and the strict-pay result
(Theorem~\ref{thm:envelope-finite}, Proposition~\ref{prop:tower},
Proposition~\ref{prop:strict}, Theorem~\ref{thm:general}); and exact machine verification
across all of it.

\noindent \textbf{Not claimed:} the discriminating-kernel question or its algorithm (Aubin;
Aubin and Catt\'e; Cardaliaguet, Quincampoix and Saint-Pierre), Farkas duality, the Isaacs
condition, minimax equality (Sion), or existence of a value under incomplete information
(Cardaliaguet and Quincampoix). The contribution is the \emph{certificate form} --- a
finitely supported, tightly bounded, exactly checkable measure witness --- and the
\emph{negative boundary result} that says precisely when that form ceases to exist.

\section{Conclusion}
\label{conclusion}

The common-action obstruction says that no single instrument serves every branch of an
information set. Its polyhedral form carries a Farkas multiplier, and a Farkas multiplier is
already a certificate: a finitely supported witness that a system of linear inequalities is
infeasible. The question this paper asked is whether that is an accident of the polyhedral
case or the shadow of something general. It is the latter, and the general object has a
shape.

\textbf{What the obstruction is dual to.} For control-affine dynamics with a compact convex
control set, the common safe-action set is empty precisely when an adversarial probability
measure on the active boundary--disturbance bundle drives the expected inward drift strictly
negative against every control (Theorem~\ref{thm:dual}). The obstruction is therefore not
merely a failure of a search to find a safe action; it is the existence of a \emph{witness} to
that failure, and the witness is a measure. This is the same movement that carries linear
programming from "the solver found nothing" to "here is a dual ray", and it carries the same
benefit: the certificate is checkable without re-running the search, and it localises the
failure on the active set.

\textbf{What governs the size of the witness.} The bound is the substantive point. The
certificate is supported on at most \(k + 1\) points, where \(k\) is the \emph{control}
dimension, and that bound is tight (Proposition~\ref{prop:sparse}) --- it is attained, not
merely approached. Neither the state dimension, nor the cardinality of the disturbance set,
nor the resolution of any discretisation appears in it. So the complexity of certifying
nonviability is governed by how many independent controls the agent has, not by how large the
world it operates in is. A companion paper reaches an analogous conclusion for a different
pipeline: there the count is the information--time product rather than the input dimension
(Abaee, 2026, computational certification). The two are the same finding in different
languages, and the recurrence is the point --- \emph{obstruction certificates are small
objects whose size is set by the decision variables, not by the uncertainty.} That is what
makes them usable.

\textbf{What the two recoveries establish.} The polyhedral case recovers the
\(\ell_{1}\)-normalised Farkas multiplier and the game case recovers the Isaacs condition
(Proposition~\ref{prop:recover}). Neither is a consistency check; each says that a classical
duality is a specialisation of this one, obtained by restricting the bundle. The measure form
is thus not a third thing alongside Farkas and Isaacs duality but the common generalisation
from which both drop out --- which is the sense in which it is the \emph{dual of the
obstruction} rather than a dual of one of its presentations.

\textbf{What the boundary result delimits.} The certificate form is not universal, and the
paper locates precisely where it stops rather than leaving that to be discovered: the convexity
boundary is identified exactly, with an explicit gap instance on the other side of it
(Proposition~\ref{prop:gap}). Knowing when a certificate cannot exist is as much a part of the
contribution as knowing what it looks like where it does, because it converts a silent failure
into a diagnosed one.

\textbf{Scope, restated once.} Claimed here is the certificate form and the negative boundary,
together with the dynamic, review-sequence and finite-sequence envelope extensions and their
exact machine verification. Not claimed are the discriminating-kernel question and its
algorithm, Farkas duality, the Isaacs condition, minimax equality, or the existence of a value
under incomplete information --- all of which are assumed, cited, or deliberately outside the
frame. The contribution is a \emph{form}: a finitely supported, tightly bounded, exactly
checkable measure witness to nonviability, and a statement of when that form ceases to exist.

\section{References}
%% ------------------------------------------------------------------
%% Added 2026-09-29 to support section \ref{priorart}. All VERIFIED against
%% publisher records (journal, volume, pages, year confirmed):
%%   Aubin 1991, Viability Theory, Birkhauser Boston
%%   Aubin & Catte 2002, Set-Valued Anal. 10, 379-416 (doi 10.1023/A:1020667819804)
%%   Cardaliaguet, Quincampoix & Saint-Pierre 1994, RAIRO Modelisation Math.
%%     Anal. Numer. (M2AN) 28(4), 441-461  -- the discriminating-kernel ALGORITHM
%%     paper; its Prop 2.3 is the convexity-of-kernel result this paper's
%%     prop:gap is the complement of.
%%   Cardaliaguet, Quincampoix & Saint-Pierre 2007, in Advances in Dynamic Game
%%     Theory, Ann. ISDG 9, 3-35, Birkhauser Boston
%%   Cardaliaguet & Quincampoix 2008, Int. Game Theory Rev. 10(1), 1-16 --
%%     value via HJI ON THE SPACE OF MEASURES under probability knowledge of the
%%     initial condition. The single closest neighbour.
%%   Sion 1958, Pacific J. Math. 8(1), 171-176 (doi 10.2140/PJM.1958.8.171)
%% NOTE ON AUTHOR NAMES: it is Aubin and CATTE (accented), not "Catte".
%%   Do not strip the accent.
%% NOTE ON A DROPPED CLAIM: an earlier roadmap note attributed to Aubin (2001)
%%   the "complement of the kernel -- Poincare's shadow". That attribution could
%%   NOT be verified and is NOT asserted here. Aubin (2001) is cited only for
%%   the verified fixed-point characterisation of kernels and capture basins.
%% ------------------------------------------------------------------
Aubin, J.-P., 1991. \emph{Viability Theory}. Systems and Control: Foundations and
Applications. Birkh\"auser Boston, Boston, MA.

Aubin, J.-P., 2001. Viability kernels and capture basins of sets under differential
inclusions. \emph{SIAM Journal on Control and Optimization} \textbf{40}(3), 853--881.

Aubin, J.-P. and Catt\'e, F., 2002. Bilateral fixed-points and algebraic properties of
viability kernels and capture basins of sets. \emph{Set-Valued Analysis} \textbf{10},
379--416.

Cardaliaguet, P. and Quincampoix, M., 2008. Deterministic differential games under
probability knowledge of initial condition. \emph{International Game Theory Review}
\textbf{10}(1), 1--16.

Cardaliaguet, P., Quincampoix, M. and Saint-Pierre, P., 1994. Some algorithms for differential
games with two players and one target. \emph{RAIRO -- Mod\'elisation Math\'ematique et Analyse
Num\'erique} \textbf{28}(4), 441--461.

Cardaliaguet, P., Quincampoix, M. and Saint-Pierre, P., 2007. Differential games through
viability theory: old and recent results. In: J\o rgensen, M., Quincampoix, M. and Vincent,
T.L. (eds), \emph{Advances in Dynamic Game Theory}, Annals of the International Society of
Dynamic Games vol.\ 9, pp.\ 3--35. Birkh\"auser Boston.

Sion, M., 1958. On general minimax theorems. \emph{Pacific Journal of Mathematics}
\textbf{8}(1), 171--176.


Abaee, A., 2026. Robust viability of the 2J3KL limit reference point
under a surplus-production map: policy scoring, expansion, and when catch
cannot help. Zenodo. \url{https://doi.org/10.5281/zenodo.22552060}

Abaee, A., 2026. An Obstruction Calculus for Viability under Incomplete Observation. Figshare preprint, version 2. \url{https://doi.org/10.6084/m9.figshare.33716593.v2}.

Abaee, A., 2026. Obstruction certificates under incomplete observation. Revised paper01 manuscript candidate (Open Problem~1); no Preprints.org DOI assigned.

Abaee, A., 2026. Exact audits of worked systems for the obstruction calculus:
counts, witnesses, and monitoring design rules (manuscript candidate).

Farkas, J., 1902. Theorie der einfachen Ungleichungen. J.\ Reine Angew.
Math. 124, 1--27.


Beiglb\"ock, M., Henry-Labord\`ere, P., Penkner, F., 2013.
Model-independent bounds for option prices --- a mass transport approach.
Finance and Stochastics 17, 477--501.

Chv\'{a}tal, V., 1983. Linear Programming. A Series of Books in the
Mathematical Sciences. W.~H. Freeman, New York.

Helly, E., 1923. \"Uber Mengen konvexer K\"orper mit gemeinschaftlichen
Punkten. J.~Deutsche Math.-Ver. 32, 175--176.

Isaacs, R., 1965. Differential Games: A Mathematical Theory with
Applications to Warfare and Pursuit, Control and Optimization. Wiley,
New York.

Abaee, A., 2026. Probabilistic sufficiency for the obstruction calculus:
belief-state safety values under partial observation (manuscript candidate).

\section*{Declarations}
\textbf{Funding.} None.\quad \textbf{Competing interests.} None.\quad
\textbf{Data availability.} No external data were used.\quad
\textbf{Code availability.} The verification script
\texttt{minimax\_dual\_certificates\_v6\_verify.py} is publicly available
at \url{https://github.com/MIKEAA2020/general-sustainability} (folder
\texttt{arena agent 1/paper rewrites/latex}); it reproduces every check family verbatim. The cross-paper certificate
exchange parser (\texttt{certificate\_exchange\_v1.py}) validates the
archived dual certificate of this paper and of the certification companion
exactly, in either direction.\quad
\textbf{AI declaration.} GLM (Z.ai), Qwen (Alibaba Cloud) and DeepSeek AI
assisted with drafting and iterative review. The author reviewed and
edited all outputs and takes responsibility for the final work.

\end{document}
